\documentclass[a4paper,12pt]{article}

\usepackage[french]{babel}
\usepackage[T1]{fontenc}
\usepackage[utf8]{inputenc}
\usepackage{amsmath, amssymb}
\usepackage{geometry}
\usepackage{graphicx}
\usepackage{array}
\usepackage{enumitem}
\usepackage{tikz}
\usepackage{pgfplots}
\pgfplotsset{compat=1.18}
\usepgfplotslibrary{statistics}

\geometry{top=2cm, bottom=2cm, left=2cm, right=2cm}

\begin{document}

\begin{center}\Large\textbf{Épreuve anticipée de mathématiques}\end{center}
\begin{center}\textbf{Voie générale : candidats suivant l'enseignement de spécialité de mathématiques.}\end{center}

\begin{center}
\textbf{Durée : 2 heures. L'usage de la calculatrice n'est pas autorisé.}
\end{center}

\section*{PREMIÈRE PARTIE : AUTOMATISMES -- QCM (6 pts)}

\textit{Pour cette première partie, aucune justification n'est demandée et une seule réponse est possible par question. Pour chaque question, reportez son numéro sur votre copie et indiquez votre réponse.}

\begin{enumerate}

\item On peut affirmer que :
\begin{itemize}
\item[A.] $\dfrac{11}{12} < \dfrac{12}{13}$
\item[B.] $\dfrac{11}{12} = \dfrac{12}{13}$
\item[C.] $\dfrac{11}{12} > \dfrac{12}{13}$
\item[D.] Il est impossible de comparer ces deux nombres.
\end{itemize}

\item Un drapeau a une aire de 2 m$^2$. Cela représente :
\begin{itemize}
\item[A.] $0,002$ cm$^2$
\item[B.] $0,2$ cm$^2$
\item[C.] $200$ cm$^2$
\item[D.] $20\,000$ cm$^2$
\end{itemize}

\item Une boîte contient 20 billes. Dans cette boîte la proportion de billes vertes est égale à $0,3$. Le nombre de billes vertes dans la boîte est égal à :
\begin{itemize}
\item[A.] $6$
\item[B.] $3$
\item[C.] $5$
\item[D.] $17$
\end{itemize}

\item Le matin à 7h00 il fait $12^\circ$C. À 10h00 la température a augmenté de $30\%$. À 10h00 il fait :
\begin{itemize}
\item[A.] $15^\circ$C
\item[B.] $16^\circ$C
\item[C.] $15,6^\circ$C
\item[D.] $42^\circ$C
\end{itemize}

\item On donne ci-contre la courbe représentative d'une fonction $g$.

\begin{center}
\begin{tikzpicture}[scale=0.8]
    \begin{axis}[
        axis lines = middle,
        xlabel = $x$,
        ylabel = $y$,
        xmin = -4, xmax = 4,
        ymin = -1, ymax = 4,
        xtick = {-3,-2,-1,0,1,2,3},
        ytick = {1,2,3},
        grid = both,
        grid style = {gray!30},
        width = 10cm,
        height = 8cm,
        samples = 100,
        smooth
    ]
        % Parabole passant par (-3;1), (0;2), (2;1)
        \addplot[thick, domain=-3.5:3.5] {-(1/9)*x^2 + 2};
        
        % Point d'ordonnée 1
        \addplot[only marks, mark=*, mark size=2] coordinates {(-3,1) (3,1)};
        
    \end{axis}
\end{tikzpicture}
\end{center}

On veut déterminer graphiquement les antécédents de 1 par $g$. On peut affirmer que :
\begin{itemize}
\item[A.] 1 n'a pas d'antécédent
\item[B.] 3 est l'antécédent de 1
\item[C.] $-3$ et $3$ sont les antécédents de 1
\item[D.] 1 a le même antécédent que $-2$
\end{itemize}

\item Parmi les nombres ci-dessous, lequel ne peut pas être une probabilité :
\begin{itemize}
\item[A.] $10^{-3}$
\item[B.] $\dfrac{20}{19}$
\item[C.] $\dfrac{1}{9}$
\item[D.] $\sqrt{2}$
\end{itemize}

\item On a résumé les effectifs d'un club de tennis de table en fonction de l'âge des adhérents dans le diagramme ci-contre.

\begin{center}
\begin{tikzpicture}
    \begin{axis}[ybar, ymin=0, ymax=25, width=10cm, height=6cm,
        xtick={1,2,3,4}, xticklabels={[0;20[, [20;40[, [40;60[, [60;80[},
        ylabel={Effectifs}, xlabel={Âge (ans)}]
        \addplot coordinates {(1,5) (2,15) (3,20) (4,15)};
    \end{axis}
\end{tikzpicture}
\end{center}

Parmi les 4 propositions ci-dessous, déterminer celle qui est assurément fausse :
\begin{itemize}
\item[A.] L'âge moyen des adhérents est de 27 ans.
\item[B.] Certains adhérents ont plus de 70 ans.
\item[C.] La moitié des adhérents a entre 20 et 40 ans.
\item[D.] Le troisième quartile de la série est égal à 54.
\end{itemize}

\item La forme factorisée de $(a+1)^2 - 3(a+1)$ est :
\begin{itemize}
\item[A.] $-(a+1)$
\item[B.] $(a+1)(-a-2)$
\item[C.] $(a+1)(a-2)$
\item[D.] $a^2 - a - 2$
\end{itemize}

\item Un train a diminué sa vitesse de $20\%$ à cause d'un incident sur la voie. Pour retrouver le niveau d'avant l'incident, le conducteur doit augmenter la vitesse de :
\begin{itemize}
\item[A.] $20\%$
\item[B.] $25\%$
\item[C.] $30\%$
\item[D.] $40\%$
\end{itemize}

\item La courbe représentative d'une fonction affine non constante est :
\begin{itemize}
\item[A.] une parabole
\item[B.] une droite horizontale
\item[C.] une droite oblique
\item[D.] une droite verticale
\end{itemize}

\item On lance un dé truqué dont les probabilités d'apparition de chaque face sont données dans le tableau ci-dessous :
\[
\begin{array}{|c|c|c|c|c|c|c|}
\hline
\text{Face} & 1 & 2 & 3 & 4 & 5 & 6 \\
\hline
\text{Probabilité} & 0,1 & 0,1 & 0,25 & 0,3 & 0,15 & 0,1 \\
\hline
\end{array}
\]
La probabilité d'obtenir un nombre inférieur ou égal à 4 est :
\begin{itemize}
\item[A.] $0,25$
\item[B.] $0,3$
\item[C.] $0,65$
\item[D.] $0,75$
\end{itemize}

\end{enumerate}

\newpage

\section*{DEUXIÈME PARTIE (14 pts)}

\subsection*{Exercice 1}

À partir du printemps 2026, Lina gère un bassin d'irrigation contenant au départ $1000$ m$^3$ d'eau. Chaque année, à cause de l'évaporation et des fuites, $20\%$ du volume d'eau est perdu. En fin d'été, Lina rajoute $40$ m$^3$ d'eau dans le bassin.

Pour tout entier naturel $n$, on note $u_n$ le volume d'eau (en m$^3$) dans le bassin au début du printemps de l'année $2026+n$. Ainsi $u_0 = 1000$.

\begin{enumerate}
\item Calculer le volume d'eau du bassin au printemps 2027 suivant ce modèle.
\item Justifier que, pour tout entier naturel $n$, $u_{n+1} = 0,8 u_n + 40$.
\item Justifier que la suite $(u_n)$ n'est pas arithmétique.
\item On considère la suite $(v_n)$ définie pour tout entier naturel $n$ par $v_n = u_n - 200$.
\begin{enumerate}
\item Montrer que la suite $(v_n)$ est géométrique de raison $0,8$.
\item En déduire l'expression de $v_n$ en fonction de $n$.
\item En déduire que, pour tout entier naturel $n$, $u_n = 800 \times 0,8^n + 200$.
\end{enumerate}
\item À l'aide d'un tableur, on a obtenu les valeurs arrondies au centième suivantes :
\[
\begin{array}{|c|c|}
\hline
n & u_n \\
\hline
5 & 409,83 \\
10 & 273,09 \\
15 & 228,36 \\
20 & 206,59 \\
25 & 201,05 \\
30 & 200,17 \\
\hline
\end{array}
\]
\begin{enumerate}
\item Conjecturer le comportement des termes de $(u_n)$ lorsque $n$ tend vers $+\infty$.
\item Interpréter ce comportement dans le contexte de l'exercice.
\end{enumerate}
\end{enumerate}

\newpage

\subsection*{Exercice 2}

Soit $f$ la fonction définie pour tout nombre réel $x \neq -\frac{1}{2}$ par
\[
f(x) = \frac{3x - 3}{2x + 1}.
\]

\begin{enumerate}
\item On admet que $f$ est dérivable sur $\mathbb{R}\setminus\left\{-\frac{1}{2}\right\}$. Montrer que pour tout $x \neq -\frac{1}{2}$,
\[
f'(x) = \frac{9}{(2x+1)^2}.
\]

\item En déduire les variations de $f$ sur son ensemble de définition.

\item On note $\mathcal{C}$ la courbe représentative de $f$.
\begin{enumerate}
\item Montrer que $\mathcal{C}$ admet exactement deux tangentes de coefficients directeurs égaux à $1$.
\item Déterminer l'équation réduite de chacune de ces tangentes.
\end{enumerate}

\item Tracer ci-dessous la courbe $\mathcal{C}$ et ses deux tangentes.

\begin{center}
\begin{tikzpicture}[scale=0.6]
    \begin{axis}[
        axis lines = middle,
        xlabel = $x$,
        ylabel = $y$,
        xmin = -5, xmax = 5,
        ymin = -5, ymax = 5,
        xtick = {-4,-3,-2,-1,0,1,2,3,4},
        ytick = {-4,-3,-2,-1,0,1,2,3,4},
        grid = both,
        grid style = {gray!30},
        width = 12cm,
        height = 12cm,
        samples = 100,
        smooth,
        clip = false
    ]
        % Asymptote verticale
        \addplot[dashed] coordinates {(-0.5,-5) (-0.5,5)};
        
        % Courbe de f
        \addplot[thick, domain=-5:-0.6] {(3*x-3)/(2*x+1)};
        \addplot[thick, domain=-0.4:5] {(3*x-3)/(2*x+1)};
        
        % Tangentes (à tracer après calcul)
        % y = x - 1 et y = x + 5
        \addplot[blue, thick, domain=-5:5] {x-1};
        \addplot[blue, thick, domain=-5:5] {x+5};
        
        % Points de tangence
        \addplot[only marks, mark=*, mark size=2, red] coordinates {(1,0) (-2,3)};
        
    \end{axis}
\end{tikzpicture}
\end{center}

\end{enumerate}

\end{document}